%------------------------------------------------------------------------------%
\section{Primitivas e Integrais Indefinidas}

%------------------------------------------------------------------------------%
\begin{definicao}{Primitiva ou Antiderivada}{def:promitiva-antiderivada}
  Uma função $F$ é denominada uma primitiva de $f$ num intervalo $I$ se
  $F'(x) = f(x)$ para todo $x \in I$.
\end{definicao}
%------------------------------------------------------------------------------%

%------------------------------------------------------------------------------%
\begin{example}~
  \solution
  A função $F(x)= \frac{1}{3}x^3$ é uma primitiva para a função $f(x)=x^2$,
  uma vez que $F'(x)=f(x)$. Note que a função $F(x)= \frac{1}{3}x^3+345$
  também é uma primitiva para $f$, o que nos leva a induzir que a
  primitiva de uma função não é única.
\end{example}
%------------------------------------------------------------------------------%

%------------------------------------------------------------------------------%
\begin{teorema}{Primitiva}{teo:primitiva}
  Se $F$ é uma primitiva de $f$ em um intervalo $I$, então a primitiva
  mais geral de $f$ em $I$ é
  \[
  F(x)+C
  \]
  onde $C$ é uma constante arbitrária.
\end{teorema}
%------------------------------------------------------------------------------%

%------------------------------------------------------------------------------%
\begin{example}
  Encontre a primitiva mais geral de cada uma das seguintes funções abaixo
  \begin{enumerate}[label=\alph*)]
    \item \(f(x)= \sin(x)\)
    \item \(f(x)= \frac{1}{x}\)
    \item \(f(x)= x^n, \quad n \neq -1\)
  \end{enumerate}

  \solution
  Usando o Teorema~\ref{teo:primitiva} em cada caso temos.
  \begin{enumerate}[label=\alph*)]
    \item
    Sabemos que se $F(x)=-\cos(x)$ é uma primitiva para $f(x)= \sin(x)$,
    uma vez que $F'(x)=\sin(x)=f(x)$. Portanto a primitiva mais geral
    para $f(x)= \sin(x)$ é $G(x)= \sin(x)+K$, onde $K$ é uma constante.
    \item
    Lembre-se que $\frac{d}{dx} (\ln(x))=\frac{1}{x}$, no intervalo
    $(0, \infty)$.
    Portanto uma primitiva mais geral no intervalo $(0,\infty)$
    para $f(x)=\frac{1}{x}$ é $G(x)=\ln(x)+c$, onde $c$ é uma constante.
    Sabemos também que $\frac{d}{dx}(\ln(\vert x\vert))= \frac{1}{x}$ em
    qualquer intervalo que não contenha  $0$, isso nos diz que
    $\ln(\vert x \vert)+ C$ é uma primitiva mais geral para
    $f(x)= \frac{1}{x}$ em cada um dos intervalos $(-\infty,0)$ e
    $(0, \infty)$.

    Em síntese, a primitiva mais geral para $f$ é
    \[
    \left\{\begin{array}{ll} \ln(x)+ C_1, \,\, se \,\, x>0\\
      \ln(-x)+ C_2 \,\, se \,\, x<0\end{array}\right.
    \]
    ou de forma mais simlificada $F(x)=\ln(\vert x \vert )+ C$,
    para todo $x \neq 0$.
    \item
    Vamor recorrer a regra da potência para encontrar uma primitiva para $x^n$.
    Para $n \neq =1$ temos que
    \[
    \frac{d}{dx} \left( \frac{x^{n+1}}{n+1}\right)
    = \frac{(n+1)x^n}{n+1}=x^n
    \]
    Logo uma primitiva mais geral para $f(x)=x^n$ é
    \[
    F(x)=\frac{x^{n+1}}{n+1}+ C
    \]
    Isso é válido para todo $n \geq 0$, uma vez que $f(x)=x^n$ está
    definida em toda a reta, nesse caso. Se $n$ for negativo
    (mas $n \neq -1$), é válido em qualquer intervalo que não contenha $x=0$.
  \end{enumerate}
\end{example}
%------------------------------------------------------------------------------%

A seguir apresentaremos, na tabela \ref{tab:primimivas}, algumas primitivas
particulares ou \textit{primitivas mais simples}, que nada mais é que a
primitiva sem a constante (ou com constate $0$). Usaremos a notação $F$ e $G$,
significando que $F$ é uma primitiva para $f$ e $G$ é uma primitiva para $g$.

%------------------------------------------------------------------------------%
\begin{table}
  \centering
  \renewcommand*{\arraystretch}{1.7}
  \begin{tabular}{ll}
    \hline
    Função                      & Primitiva particular    \\ \hline\hline
    \(cf(x)\)                   & \(cF(x) \)              \\
    \(f(x)+g(x)\)               & \(F(x)+G(x)\)           \\
    \(x^n \,\, x \neq -1\)      & \(\frac{x^{n+1}}{n+1}\) \\
    \(\frac{1}{x}\)             & \(\ln(\vert x \vert)\)  \\
    \(e^x\)                     & \(e^x\)                 \\
    \( \cos(x)\)                & \(\sin(x)\)             \\
    \(\sec^2(x)\)               & \(\tan(x)\)             \\
    \(\sec(x) \tan(x)\)         & \(\sec(x)\)             \\
    \(\frac{1}{\sqrt{1-x^2}}\)  & \(\arcsin(x)\)          \\
    \(-\frac{1}{\sqrt{1-x^2}}\) & \(\arccos(x)\)          \\
    \(\frac{1}{1+x^2}\)         & \(\arctan(x)\)          \\
    \(\sin(x) \)                & \(-\cos(x)\)            \\ \hline
  \end{tabular}
  \caption{Tabela de fórmulas para primitivas}
  \label{tab:primimivas}
\end{table}
%------------------------------------------------------------------------------%

%------------------------------------------------------------------------------%
\begin{example}
  Encontre todas as funções $g$ tais que
  \[
  g'(x)= 4 \sin(x)+ \frac{2x^5- \sqrt{x}}{x}
  \]

  \solution
  Nesse exercício queremos encontrar uma primitiva para $g'$.
  Podemos reescrever $g'$ como
  \[
  g'(x)
  = 4 \sin(x) + \frac{2x^5- \sqrt{x}}{x}
  = 4 \sin(x) + 2x^4 -x^{-\sfrac{1}{2}}
  \]
  Assim, usando a Tabela~\ref{tab:primitivas} e o
  Teorema~\ref{teo:primitivas}, temos que
  \begin{align*}
    g(x) & = 4(-\cos(x))+ 2 \frac{x^5}{5}-\frac{x^{\frac{1}{2}}}{\frac{1}{2}}+C \\
    & = -4 \cos(x)+ \frac{2}{5}x^5-2 \sqrt{5}+C
  \end{align*}
\end{example}
%------------------------------------------------------------------------------%

%------------------------------------------------------------------------------%
\begin{example}
  Encontre $f$ se $f'(x)= e^x+20(1+x^2)^{-1}$ e $f(0)=-2$.

  \solution
  Inicialmente vamos encontrar uma família de funções $f$ que são primitivas de
  $f'$, depois usamos a condição $f(0)=-2$ para encontrarmos a primitiva
  pedida.

  A primitiva geral de
  \[
  f'(x)= e^x+20 \frac{1}{1+x^2}
  \]
  é
  \[
  f(x)=e^x+20 \arctan(x)+C
  \]
  Para encontrarmos a constante $C$, usamos o fato de que $f(0)=-2$:
  \[
  f(0)=e^0+20 \arctan(0)+C=-2
  \]
  Com isso, $C=-3$ e a solução particular é
  \[
  f(x)=e^x+20 \arctan(x)-3
  \]
\end{example}
%------------------------------------------------------------------------------%
